\documentclass[11pt]{article}

\usepackage[margin=1in]{geometry}
\usepackage{booktabs}
\usepackage{amsmath}
\usepackage{hyperref}
\usepackage{xcolor}
\usepackage{longtable}
\usepackage{parskip}
\usepackage{caption}
\usepackage{fancyhdr}
\usepackage{titlesec}

\titleformat{\section}{\normalfont\Large\bfseries}{\thesection}{1em}{}
\titleformat{\subsection}{\normalfont\large\bfseries}{\thesubsection}{1em}{}

\pagestyle{fancy}
\fancyhf{}
\fancyhead[L]{Prediction Accuracy Report}
\fancyhead[R]{\thepage}
\renewcommand{\headrulewidth}{0.4pt}

\hypersetup{
    colorlinks=true,
    linkcolor=blue,
    urlcolor=blue,
    citecolor=blue,
    pdftitle={Prediction Algorithm Accuracy Report},
    pdfauthor={Sports Analytics Project}
}

\newcommand{\code}[1]{\allowbreak\texttt{\small #1}\allowbreak}
\usepackage[htt]{hyphenat}

\title{\textbf{Prediction Algorithm Accuracy Report} \\ \large Backtests and Permutation Tests Against Live Season Data}
\author{Sports Analytics Project}
\date{September 2026}

\begin{document}

\maketitle

\begin{abstract}
\noindent This report backtests all four prediction algorithms currently running in the sports analytics
platform against the project's live database: 1,321 completed NBA games (2025--26 season, regular season and
playoffs), 530 players, and 33,590 individual player-game boxscore rows. For each algorithm we report a
walk-forward accuracy metric (predictions built only from information available strictly before the event
being predicted, matching how each model is actually used in production), a comparison against a naive
baseline, and a permutation test establishing whether the observed accuracy reflects real signal rather than
chance. All four algorithms beat their naive baselines by a statistically significant margin
($p < 0.001$ in every case), but the size of that edge varies enormously — from decisive (win probability) to
marginal (point margin). We close with a per-algorithm assessment of where accuracy plausibly tops out given
more data and modeling effort, and what is worth prioritizing next.
\end{abstract}

\tableofcontents

\section{Scope and Method}

\subsection{What is being evaluated}

The codebase has four distinct prediction algorithms, spread across a NestJS API and two standalone Python
services:

\begin{enumerate}
    \item \textbf{Elo rating} (\code{apps/predictor/elo.py}) — predicts home team win probability for a game.
    \item \textbf{Four Factors regression} (\code{apps/predictor/four\_factors.py}) — predicts the home-minus-away
    point margin for a game.
    \item \textbf{Recency-weighted fantasy points} (\code{apps/optimizer/predict.py}) — predicts a player's
    DraftKings-style fantasy point output for their next game, feeding the lineup optimizer.
    \item \textbf{Recency-weighted scorer points} (\code{apps/api/src/games/game-detail.service.ts}) — predicts a
    player's raw point total for a specific upcoming matchup, feeding the game detail page's ``predicted top
    scorers'' list.
\end{enumerate}

\subsection{Data}

All backtests ran against the project's live Supabase-hosted Postgres database, not synthetic or seed data:
1,321 completed games spanning 2025-10-21 through 2026-06-13, 30 teams, 530 players, and 33,590
\code{PlayerGameStat} rows. This is a genuine full-season sample, large enough for the accuracy numbers below
to be meaningful rather than an artifact of a handful of games.

\subsection{Walk-forward evaluation, not in-sample evaluation}

A prediction is only informative about a model's real-world usefulness if it was built using nothing the
model could not have known at the time. All four production algorithms already enforce this at the level of
an individual prediction — every one of them is documented and tested (via \code{test\_four\_factors.py} and
\code{games.e2e-spec.ts}) to never let a game's own outcome, or a later game's outcome, leak into that game's
prediction.

One subtlety required extra care. \code{four\_factors.py}'s regression path
(\S\ref{sec:four-factors}) fits a single ordinary-least-squares model once, using every game in the season,
and then applies that one fitted model to predict every game — including games early in the season, whose
prediction is therefore informed (through the fitted regression weights, not through any single feature
value) by results that happened after them. This is leak-free per game but not a true walk-forward
backtest. We report both numbers: the \emph{as-shipped} evaluation (exactly what \code{predict\_games.py}
currently produces) and a \emph{true walk-forward} evaluation that periodically refits the regression using
only strictly-past games. The gap between them tells us how much the as-shipped number is inflated by
full-season fitting.

The other three algorithms (Elo, and both recency-weighted point predictors) are walk-forward by
construction — Elo's rating update is inherently sequential, and the recency-weighted average at
game $i$ only ever consumes games $0, \ldots, i-1$ for a given player or team.

\subsection{Metrics}

\begin{itemize}
    \item \textbf{Win probability (Elo)}: Brier score (mean squared error between predicted probability and
    the 0/1 outcome; 0 = perfect, 0.25 = uninformative coin flip), log loss, and thresholded accuracy
    (predict the team with $P > 0.5$).
    \item \textbf{Margin (Four Factors)}: mean absolute error (MAE) and root-mean-squared error (RMSE) in
    points.
    \item \textbf{Player point predictors}: MAE and RMSE in points (or fantasy points), each compared against
    an unweighted running mean of the same prior games — this isolates whether \emph{recency weighting
    specifically} adds value over simply averaging a player's history.
\end{itemize}

\subsection{Permutation tests}

For each algorithm we constructed a null distribution appropriate to how that specific model consumes its
data, then computed a p-value as the fraction of null-model runs that matched or beat the real model's
observed accuracy:

\begin{itemize}
    \item \textbf{Elo}: re-ran the entire sequential Elo forward pass 1,000 times on the real schedule, but
    with each game's outcome replaced by an independent coin flip. This tests whether Elo's Brier score
    reflects real team-strength signal, not just structural artifacts of the home-court constant or schedule.
    \item \textbf{Four Factors}: shuffled the realized-margin labels across games 1,000 times (keeping the
    leak-free feature matrix fixed) and refit the regression on each shuffled copy. This tests whether the
    Four Factors $\to$ margin relationship is real signal or an artifact of fitting three free parameters
    against noise.
    \item \textbf{Recency-weighted predictors}: shuffled each player's own game order 500 times (destroying
    any genuine hot/cold-streak or role-change structure while preserving each player's overall scoring
    level) and recomputed the recency-weighted model's edge over the naive running mean on the shuffled
    sequences. This tests whether the recency-weighting edge is really about recency, or would appear even
    against a randomly ordered history.
\end{itemize}

\section{Results}

\subsection{Elo win probability}
\label{sec:elo}

\begin{table}[h]
\centering
\begin{tabular}{lr}
\toprule
Metric & Value \\
\midrule
Games evaluated & 1,321 \\
Brier score & 0.2146 \\
Log loss & 0.6179 \\
Accuracy (threshold 0.5) & 65.2\% \\
Naive baseline (``home team always wins'') & 55.5\% \\
Mean predicted $P(\text{home win})$ & 0.601 \\
\bottomrule
\end{tabular}
\caption{Elo win-probability backtest, as-shipped (walk-forward by construction).}
\end{table}

A Brier score of 0.2146 against the 0.25 coin-flip baseline, and 65.2\% thresholded accuracy against a 55.5\%
naive-home-favorite baseline, is a real and useful edge for a model that uses only game scores — no injuries,
rest, travel, or player-level information at all.

\textbf{Calibration.} Binning predictions into deciles and comparing mean predicted probability to mean
actual win rate in each bin shows good calibration in the upper range (0.6--0.9 predicted probability bins
track actual outcomes closely) but a moderate overconfidence problem in the 0.2--0.5 range, where the model's
predicted probabilities run consistently higher than the observed win rate (e.g. the 0.4--0.5 bin predicts
0.454 on average but the true rate in that bin was 0.359).

\begin{table}[h]
\centering
\small
\begin{tabular}{lrrr}
\toprule
Predicted range & $n$ & Mean predicted & Mean actual \\
\midrule
0.1--0.2 & 5 & 0.173 & 0.000 \\
0.2--0.3 & 27 & 0.257 & 0.037 \\
0.3--0.4 & 87 & 0.357 & 0.230 \\
0.4--0.5 & 181 & 0.454 & 0.359 \\
0.5--0.6 & 308 & 0.554 & 0.451 \\
0.6--0.7 & 384 & 0.644 & 0.646 \\
0.7--0.8 & 230 & 0.746 & 0.765 \\
0.8--0.9 & 92 & 0.839 & 0.848 \\
0.9--1.0 & 7 & 0.920 & 0.857 \\
\bottomrule
\end{tabular}
\caption{Elo calibration, 10 equal-width probability bins.}
\end{table}

The consistent pattern across the 0.2--0.6 range (predicted probability higher than actual outcome) suggests
the home-court-advantage constant (\code{HOME\_COURT\_ADVANTAGE\_ELO}, set to 75 and borrowed from
FiveThirtyEight's published range rather than fit to this league/season) may be somewhat too large for the
current data, or that \code{K\_FACTOR} (set to 20) is letting ratings drift further from 1500 than the data
actually supports before they've had enough games to justify it.

\textbf{Permutation test.} Across 1,000 randomized-outcome reruns of the same schedule, the null Brier score
distribution had mean 0.2668 and standard deviation 0.0034; the observed Brier score of 0.2146 was better
than every single one of the 1,000 null runs ($p < 0.001$). Elo's edge is real signal, not an artifact of the
home-court constant or schedule structure alone.

\subsection{Four Factors margin}
\label{sec:four-factors}

\begin{table}[h]
\centering
\begin{tabular}{lrr}
\toprule
Metric & As-shipped (full-season fit) & True walk-forward \\
\midrule
$n$ & 1,305 & 1,239 \\
MAE (points) & 12.52 & 12.65 \\
RMSE (points) & 15.80 & 15.94 \\
\bottomrule
\end{tabular}
\caption{Four Factors margin backtest. Walk-forward refits the regression every 82 games using only
strictly-past data; the small gap (+0.13 MAE) shows the as-shipped full-season fit is only mildly optimistic.}
\end{table}

\begin{table}[h]
\centering
\begin{tabular}{lr}
\toprule
Model & MAE (points) \\
\midrule
Four Factors regression (as-shipped) & 12.52 \\
Naive baseline (predict leaguewide average home margin) & 13.36 \\
\bottomrule
\end{tabular}
\caption{Four Factors vs. the simplest possible baseline.}
\end{table}

This is the weakest result of the four. The fitted regression beats the naive ``always predict the
leaguewide average home margin'' baseline by only 0.84 points of MAE — about a 6\% improvement. The
permutation test confirms this small edge is still statistically real (observed MAE 12.52 vs. a null
distribution with mean 13.34 and standard deviation only 0.02 — every one of 1,000 label-shuffled refits
performed worse, $p < 0.001$), so the Four Factors $\to$ margin relationship is not noise. But an
average absolute error of 12.5 points on an NBA game (where the typical margin is well within a few
possessions of a coin flip) means this number is not yet precise enough to be useful for anything beyond
directional (favored-team) guidance, which the Elo win-probability output already provides more reliably.

It is worth noting that the regression path \emph{is} live in production against the current dataset (1,305
of 1,321 games used the fitted regression, comfortably above the \code{MINIMUM\_GAMES\_FOR\_REGRESSION = 30}
threshold) — this is a change from the assumption in the code's own comments, which describe the heuristic
fallback as what the project's ``seed data'' would exercise. With a full season loaded, the regression is
what actually runs.

\subsection{Recency-weighted fantasy points (optimizer)}
\label{sec:fantasy}

\begin{table}[h]
\centering
\begin{tabular}{lrr}
\toprule
Metric & Recency-weighted & Naive running mean \\
\midrule
$n$ (player-game predictions) & \multicolumn{2}{c}{30,447} \\
MAE (fantasy points) & 7.74 & 8.25 \\
RMSE (fantasy points) & 10.39 & 10.96 \\
\bottomrule
\end{tabular}
\caption{Fantasy-point prediction, recency-weighted average (\code{RECENCY\_DECAY = 0.8}, unbounded window)
vs. an unweighted running mean of the same prior games.}
\end{table}

Recency weighting improves MAE by 0.51 fantasy points over a plain running mean of the same games — a modest
but consistent edge (about 6\%). The permutation test, which reshuffles each player's own game order 500
times, shows this edge disappears (and reverses to slightly negative, mean $-0.277$) once the real
chronological structure is destroyed, with every one of 500 shuffles performing worse than the observed
edge ($p < 0.001$). This confirms the model is capturing genuine within-season trend (role changes, form,
minutes changes) rather than just being a differently-weighted average that happens to fit better.

\subsection{Recency-weighted scorer points (game detail page)}
\label{sec:scorers}

\begin{table}[h]
\centering
\small
\begin{tabular}{lrr}
\toprule
Metric & \shortstack{Recency-weighted\\(10-game window)} & \shortstack{Naive running mean\\(all prior games)} \\
\midrule
$n$ (player-game predictions) & \multicolumn{2}{c}{30,447} \\
MAE (points) & 4.45 & 4.64 \\
RMSE (points) & 6.12 & 6.28 \\
\bottomrule
\end{tabular}
\caption{Raw scoring-point prediction, capped 10-game recency window vs. unweighted running mean.}
\end{table}

Same pattern as \S\ref{sec:fantasy}, smaller in absolute terms since raw points have a smaller range than
the fantasy-point composite: a 0.195-point MAE edge over the naive mean, statistically real
($p < 0.001$ against a permuted-order null with mean edge $-0.199$).

\section{Cross-Cutting Observations}

\begin{itemize}
    \item \textbf{Every algorithm beats its naive baseline by a statistically significant margin.} None of the
    four models' apparent accuracy is an illusion of overfitting or chance — all four permutation tests
    reject the null at $p < 0.001$.
    \item \textbf{The size of the edge varies by an order of magnitude.} Elo's edge over picking the home
    favorite (65.2\% vs 55.5\% accuracy) is large and immediately useful. Four Factors' edge over predicting
    the league-average margin (6\% MAE improvement) is real but small enough that the current margin number
    is not precise enough to be a standalone product feature yet.
    \item \textbf{None of the four models use any player-availability, injury, rest, or travel signal.} All
    four are pure box-score/outcome-history models. This is very likely the single largest source of
    remaining error across the board, and it is untouched by any of the modeling choices evaluated here.
    \item \textbf{The Four Factors regression is already running on real (non-heuristic) fitted weights},
    contrary to what the code comments describe as the expected state with ``seed data'' — worth updating
    those comments so they don't mislead a future reader about which code path is actually live.
\end{itemize}

\section{Where Each Model Can Realistically Go}

\subsection{Elo win probability}

Current state: a solid, well-calibrated baseline with one identifiable weakness (overconfidence in the
0.3--0.5 predicted range). Two specific, low-effort next steps:

\begin{enumerate}
    \item \textbf{Margin-of-victory scaling on the $K$-factor} — already flagged as a documented future step
    in \code{elo.py}'s own comments. Blowout wins currently move a team's rating by exactly as much as a
    one-point win, which is known (from the wider Elo/sports-analytics literature) to leave real information
    on the table. This is a small, well-understood change.
    \item \textbf{Re-derive \code{HOME\_COURT\_ADVANTAGE\_ELO} and \code{K\_FACTOR} from this project's own
    data} rather than literature defaults, given the calibration bias observed in \S\ref{sec:elo}. With
    1,321 games now available, a simple grid search over a held-out validation split is enough to do this
    properly — there was not enough data to justify this before, there is now.
\end{enumerate}

Realistic ceiling without adding new input data: Brier score in the 0.19--0.20 range (vs. today's 0.2146),
comparable to published results for pure-Elo NBA models. Getting meaningfully below that requires the
injury/rest/lineup signal described below, which is a different category of change, not a tuning exercise.

\subsection{Four Factors margin}

This is the model furthest from being a reliable product feature today. Two factors are plausible sources
of the weak signal:

\begin{enumerate}
    \item \textbf{Two of Oliver's four factors are missing by design} (offensive rebound rate and the
    defensive-side factors — the code's own docstring explains why, correctly, given the current schema's
    inability to distinguish offensive from defensive rebounds). Adding an offensive/defensive rebound split
    to \code{PlayerGameStat} ingestion would let a true four-factor model be fit instead of the current
    three-factor offense-only version, which is the most direct way to close some of this gap.
    \item \textbf{No opponent-adjustment.} The current features are each team's own raw Four Factors
    averages, with no accounting for the strength of the opponents that produced those averages — a team
    that has played a soft schedule looks better than it is. An opponent-adjusted (strength-of-schedule
    weighted) version of the same three factors is a natural next iteration and does not require new data,
    only a different aggregation of what already exists.
\end{enumerate}

Realistic ceiling: point-margin prediction in the broader sports-analytics literature (with far richer
inputs — play-by-play, lineup data, injury reports) typically sits around 10--11 points MAE for NBA games,
which is itself close to the theoretical noise floor of single-game outcomes. Given that, this project
should treat \emph{win probability} (Elo) as the primary product signal and treat the margin number as a
secondary, lower-confidence figure until the opponent-adjustment and rebound-split changes above are made —
promising it as equally reliable today would overstate what it can do.

\subsection{Recency-weighted player point predictors (both)}

Both player-level models show the same profile: a real but modest edge from recency weighting alone, with
plenty of room to grow by adding features neither model currently has access to:

\begin{enumerate}
    \item \textbf{Opponent defensive strength.} Neither model adjusts for who the player is about to face —
    a player's predicted points are currently identical whether they are playing the league's best or worst
    defense. This is likely the single highest-value addition for both models, and unlike some of the other
    recommendations, doesn't require new data ingestion — opponent identity is already available on every
    upcoming \code{Game} row.
    \item \textbf{Minutes/role trend.} Both models predict scoring output directly, with no separate signal
    for playing time. A player returning from injury with reduced minutes, or a player who has just moved
    into a starting role, will be systematically mispredicted by a pure recency-weighted average of their
    box score until enough new games accumulate to catch up. A separate minutes-trend feature (or predicting
    per-minute rate and multiplying by a minutes forecast) is the standard fix.
    \item \textbf{Home/away split.} Neither model currently distinguishes a player's home vs. away scoring
    history, despite this being one of the more well-established effects in basketball analytics.
\end{enumerate}

Realistic ceiling: the specific numbers above (MAE around 7.7 fantasy points / 4.5 raw points) are unlikely
to improve dramatically from parameter tuning alone (e.g. adjusting \code{RECENCY\_DECAY} or the game-window
size) — the permutation tests show recency weighting's contribution over a plain average is real but modest,
so most of the room for improvement is in adding the features listed above, not in re-tuning the existing
one-feature model.

\section{Recommendations, in Priority Order}

\begin{enumerate}
    \item \textbf{Add an accuracy-tracking job.} None of these four algorithms currently have any ongoing
    accuracy monitoring — this report is, as far as this codebase's test suite goes, the first time any of
    them have been evaluated against realized outcomes rather than just checked for data-leakage correctness.
    A lightweight scheduled job that recomputes Brier score / MAE against newly completed games (reusing the
    scripts built for this report) would catch model drift going forward rather than requiring another
    one-off audit.
    \item \textbf{Add opponent-adjustment to the two player-point predictors} (\S\ref{sec:fantasy},
    \S\ref{sec:scorers}) — highest expected value per unit of engineering effort, since the needed data
    (opponent team identity) is already in the schema.
    \item \textbf{Re-tune Elo's \code{HOME\_COURT\_ADVANTAGE\_ELO} and \code{K\_FACTOR}} against this
    project's own 1,321-game sample instead of literature defaults, to address the calibration bias found in
    \S\ref{sec:elo}.
    \item \textbf{Treat the Four Factors margin number as secondary to Elo's win probability} in any
    user-facing framing until the opponent-adjustment and offensive/defensive rebound split
    (\S\ref{sec:four-factors}) are implemented — today's margin prediction is only marginally better than
    guessing the league average.
    \item \textbf{Update code comments in \code{four\_factors.py}} that describe the regression path as
    unreachable with current data — it is reachable and already running.
\end{enumerate}

\section*{Appendix: Reproducing These Results}

All four backtest scripts import the real production modules (\code{elo.py}, \code{four\_factors.py},
\code{predict.py}) directly rather than reimplementing their logic, so results reflect the code as actually
shipped, not a paraphrase of it. Each script connects to the project's database (read-only for the backtest
itself), computes the walk-forward metrics and permutation test described above, and prints a plain-text
summary. Random seeds are fixed (42/43) for reproducibility of the permutation null distributions.

\end{document}
